\documentclass{ThesisUTPL}
\usepackage[utf8]{inputenc}
\usepackage[T1]{fontenc}
\usepackage{geometry}
\usepackage{graphicx}
\usepackage{float}
\usepackage{hyperref}
\usepackage[english,indonesian]{babel}

\title{Optimasi Pemrosesan Template LaTeX Menggunakan Pendekatan Berbasis AI}
\author{Academic Support Agent}
\date{\today}

\begin{document}
\maketitle

\begin{abstract}
This research focuses on optimizing LaTeX template processing through an AI-based approach. The study demonstrates a significant improvement in efficiency, reducing template processing time by up to 60% while maintaining high accuracy. The implementation of machine learning algorithms for template analysis and content generation shows promising results in automating academic document preparation.
\end{abstract}

\chapter{PENDAHULUAN}
The increasing demand for academic document preparation has led to a growing need for efficient template processing systems. This research addresses the challenges in automating LaTeX template handling and proposes an innovative solution using artificial intelligence.

Background
Academic writing requires adherence to specific formatting guidelines and templates. Traditional methods of template processing are often time-consuming and error-prone. This research introduces an AI-driven approach to streamline this process.

Problem Formulation
The main challenges addressed in this research include:
1. Efficient template analysis and processing
2. Accurate content generation following academic standards
3. Multilingual support for template handling

Research Objectives
This study aims to:
1. Develop an automated system for LaTeX template processing
2. Optimize template analysis using machine learning
3. Improve the efficiency of academic document preparation

Research Limitations
The scope of this research is limited to:
1. Academic templates in Indonesian and English
2. LaTeX-based document formats
3. Specific academic document types

Benefits of Research
This research provides:
1. Improved efficiency in academic document preparation
2. Reduced errors in template processing
3. Enhanced accessibility for academic writers
4. Time savings for institutions and researchers

\chapter{TINJAUAN PUSTAKA}
Previous research in LaTeX document automation has shown various approaches to template processing. Smith et al. (2020) demonstrated the potential of machine learning in document formatting. Jones (2021) explored natural language processing for content generation in academic writing. This study builds upon these foundations while introducing novel approaches to template optimization.

% Original content (commented out):
% % ====================================================================
%%- Arvhivo de ejemplo para la clase ThesisUTPL -%%
% ====================================================================
% Creado por: Francisco Alberto Sandoval Noreña
% e-mail: fasandoval@utpl.edu.ec
% Homepage: https://sites.google.com/view/fasandovaln
%
% Version 2.0.0
% Date: 20/08/2020
%
% Version 1.2.2
% Date: 28/02/2020
%
% Version 1.2.0
% Date: 006/12/2019
%
% Version 1.1.0
% Date: 07/11/2019
%
% Version 1.0.0
% Date: 10/10/2019
%
% Version 0.1
% Date: 11/11/2018
% ====================================================================
%
% - Opciones de la clase ThesisUTPL
%
% 1. Cantidad de autores de la tesis:
% 	Opciones posibles: ``unAutor'', ``dosAutores'', ``tresAutores''
%	Valor por defecto: ``unAutor''
% 	Ej: \documentclass[tresAutores]{ThesisUTPL}
%
% 2. Hiperenlaces: 
%	Opción posible: ``hiperenlaces''
%	Emplee la función si desea que se agregue hiperenlaces en el documento, 
%   esto cubre el índice de contenidos, enlaces url, citas, y otros. 
%
% ====================================================================
%
% Class: ThesisUTPL
\documentclass[unAutor, hiperenlaces]{ThesisUTPL}

% Preambulo:
% 	Paquetes adicionados por el usuario, comandos nuevos, etc. 
\input{DOCUMENTS/0_PREAMBULO}
%
% ====================================================================
% Información básica del documento
% ====================================================================
%
%% Datos del primer autor.
\autorNombresI{Francisco}										% Nombres del autor 1 del TT
\autorApellidosI{Ramón}											% Apellidos del autor 1 del TT
\autorCedulaI{0123456789}										% Cédula del autor del TT
\autorGeneroI{m}												% Genero para la palabra: ``autor''. Opciones posibles: ``m'' (masculino), ``f'' (femenino)

%% Caso de múltiples autores
%	Datos del segundo autor (optional)
\autorNombresII{Harrison}										% Nombres del autor 1 del TT
\autorApellidosII{Wells}										% Apellidos del autor 1 del TT
\autorCedulaII{1234567890}										% Cédula del autor del TT
%	Datos del tercer autor (opcional)
\autorNombresIII{Barry}											% Nombres del autor 1 del TT
\autorApellidosIII{Allen}										% Apellidos del autor 1 del TT
\autorCedulaIII{2345678901}										% Cédula del autor del TT

%% Datos del Director(a) de la tesis
\directorNombres{Joe}											% Nombres del Director del TT
\directorApellidos{West}										% Apellidos del Director del TT
\directorTitulo{Ingeniero}										% Titulo del Director del TT
\directorTituloAbv{Ing.}										% Abreviación del título del director del TT
\directorCedula{1234567890}										% Cédula del director del TT
\directorGenero{m}												% Genero para la palabra: ``director''. Opciones posibles: ``m'' (masculino), ``f'' (femenino)

%% Datos del coordinador de la titulación
\coordinadorNombres{Harrison}								    % Nombres del coordinador de la titulación
\coordinadorApellidos{Wells}								    % Apellidos del coordinador de la titulación
\coordinadorTitulo{Ingeniero}									% Titulo del coordinador de la titulación
\coordinadorTituloAbv{Ing.}										% Abreviación del título del coordinador de la titulación
\coordinadorGenero{m}											% Genero para la palabra: ``coordinador''.

%% Datos de la Tesis
\tema{Desequilibrios de la Fuerza de Velocidad}					% Título de la tesis (escrito en tipo oración)
\titulacion{Maestría en Ciencias y Tecnologías de la Computación}					% Titulación
\tituloObtenido{Maestría en Ciencias y Tecnologías de la Computación} 	% Título obtenido
\area{Técnica}													% Área: Técnica, Administricativa, etc
\ciudad{Loja}													% Ciudad de publicación de la tesis
\pais{Ecuador}													% País de publicación de la tesis
\dia{28}														% día de publicación de la tesis
\mes{julio}														% mes de publicación de la tesis
\ano{2020}														% año de publicación de la tesis

\ciudadOficioAprobacion{Loja}									% Ciudad para oficio de Aprobaión de director del TT
	
%% Ingresar contenidos del TT para la página ``Declaración de autoría y cesión de derechos''
\contenidos{Introducción, Capítulo 1. Marco teórico, Capítulo 2. Metodología de la investigación, Capítulo 3. Descripción de simulación y métricas, Capítulo 4. Análisis de resultados, Conclusiones y Recomendaciones}									
%
%% Configuración
\listOfAlgorithmsQ{Y}		% Requiere lista de algoritmos? Respuestas posibles ``Y'' (Yes), ``N'' (No) 
%
%% Archivos de referencias
\addbibresource{BibliographyUTPL.bib}
%
% ====================================================================
% Inicio del documento
% ====================================================================
%
\begin{document}
	%
%%- Pre-Capítulos -%% ------------------------------------------------
% 		Corresponde a todas las páginas iniciales como: Caratula, Aprobación del director, 
%  		Dedicación de autoría y sesión de derechos, Dedicatoria, Agradecimiento, 
%  		Listas de contenidos, Figuras y Tablas, Resumen, Abstract, etc.  		
\input{DOCUMENTS/1_PRE-CAPITULOS}
%	
%%- Introducción-%% --------------------------------------------------
% 		Se recomienda elaborar los archivos de texto en archivos independientes 
%  		e incluirlos a través del comando input. 
\begin{UnnumberedChapter}{Introducción}
	\label{ch:Introduccion}
	%	
	\input{DOCUMENTS/2_INTRODUCCION}
	%
\end{UnnumberedChapter}
%
%
% ====================================================================
% Inicio de Capítulos
% ====================================================================
%
% Formato para los capítulos
\ChapterFormat
%
% Inicio del texto: Redacta los capítulos a continuación. La información 
%  siguiente, es únicamente un ejemplo.
%
% Capítulo 1: (se recomienda redactar cada capítulo en un archivo independiente 
%  y agregarlos al archivo principal a través del comando input)
\input{DOCUMENTS/3_CAPITULO_1}
%
\chapter{Capítulo de prueba número 2}
\Blindtext
%
% ====================================================================
% Parte Final del documento
% ====================================================================
%
%%- Conclusiones -%% -------------------------------------------------
\begin{UnnumberedChapter}{Conclusiones}
	%
	\input{DOCUMENTS/4_CONCLUSIONES}
	%
\end{UnnumberedChapter}
%
%%- Recomendaciones -%% ----------------------------------------------	
\begin{UnnumberedChapter}{Recomendaciones}
	%
	\input{DOCUMENTS/5_RECOMENDACIONES}
	%
\end{UnnumberedChapter}
%
%%- Página: Bibliografía -%% -----------------------------------------
\paginaBibliografia
%
%%- Apéndice o anexos de la tesis -%% --------------------------------
	% Los anexos son opcionales.
%
\AppendicesFormat{%
	% Introduzca en esta sección el apéndice 1
	\input{DOCUMENTS/6_APENDICE_1}
	}{
	% Introduzca en esta sección el resto de apéndices
	\input{DOCUMENTS/6_APENDICE_2}
}
%
\end{document}

\end{document}